% !TeX program = pdflatex
\documentclass[12pt,a4paper]{article}

\newcommand{\notelanguage}{finnish}
% Shared preamble for course notes and exercise sets.

\usepackage[T1]{fontenc}
\usepackage{lmodern}
\usepackage[english]{babel}

\usepackage[a4paper,margin=30mm]{geometry}
\usepackage{microtype}
\usepackage{graphicx}
\usepackage{enumitem}

\usepackage{amsmath}
\usepackage{amssymb}
\usepackage{amsthm}
\usepackage{mathtools}
\usepackage{aliascnt}

\usepackage[hidelinks,pdfusetitle]{hyperref}

\numberwithin{equation}{section}
\setlist{itemsep=0.25em,topsep=0.5em}

\theoremstyle{plain}
\newtheorem{theorem}{Theorem}[section]
\newaliascnt{lemma}{theorem}
\newtheorem{lemma}[lemma]{Lemma}
\aliascntresetthe{lemma}
\newaliascnt{proposition}{theorem}
\newtheorem{proposition}[proposition]{Proposition}
\aliascntresetthe{proposition}
\newaliascnt{corollary}{theorem}
\newtheorem{corollary}[corollary]{Corollary}
\aliascntresetthe{corollary}

\theoremstyle{definition}
\newaliascnt{definition}{theorem}
\newtheorem{definition}[definition]{Definition}
\aliascntresetthe{definition}
\newaliascnt{example}{theorem}
\newtheorem{example}[example]{Example}
\aliascntresetthe{example}
\newaliascnt{problem}{theorem}
\newtheorem{problem}[problem]{Problem}
\aliascntresetthe{problem}

\theoremstyle{remark}
\newaliascnt{remark}{theorem}
\newtheorem{remark}[remark]{Remark}
\aliascntresetthe{remark}

\usepackage[nameinlink,noabbrev]{cleveref}

\crefname{theorem}{theorem}{theorems}
\crefname{lemma}{lemma}{lemmas}
\crefname{proposition}{proposition}{propositions}
\crefname{corollary}{corollary}{corollaries}
\crefname{definition}{definition}{definitions}
\crefname{example}{example}{examples}
\crefname{problem}{problem}{problems}
\crefname{remark}{remark}{remarks}

\DeclareMathOperator{\im}{im}
\DeclareMathOperator{\rank}{rank}
\DeclarePairedDelimiter{\abs}{\lvert}{\rvert}
\DeclarePairedDelimiter{\norm}{\lVert}{\rVert}
\DeclarePairedDelimiter{\set}{\lbrace}{\rbrace}


\usepackage{tikz}

\title{Ryhmätyö 2 ratkaisut}
\author{}
\date{}

\begin{document}

\maketitle

\noindent

\begin{exercise}
  \textbf{Määritelmä.} Lukujono $(x_n)$ on \emph{ylhäältä rajoitettu}, jos on
  olemassa reaaliluku $M$, jolle $x_n \leq M$ kaikilla $n$.
\begin{enumerate}[label=\alph*), nosep]
  \item Määrittele vastaava käsite \emph{alhaalta rajoitettu}.
  \item Anna esimerkki lukujonosta, joka on ylhäältä rajoitettu mutta ei ole
    alhaalta rajoitettu.
  \item Anna esimerkki lukujonosta, joka ei ole ylhäältä eikä alhaalta
    rajoitettu.
  \item Lukujono on \emph{rajoitettu}, jos se on ylhäältä ja alhaalta
    rajoitettu. Onko tämä ekvivalentti kirjan antaman määritelmän kanssa
    (ennen Lausetta 2.2.7)?
\end{enumerate}
\end{exercise}

\begin{solution}
\leavevmode\par
\begin{enumerate}[label=\alph*), nosep]
\item \textbf{Määritelmä.} Lukujono \((x_n)\) on \emph{alhaalta rajoutettu},
  jos on olemassa reaaliluku \(M\), jolle \(M \leq x_n\) kaikilla \(n \in \NN\).
\item Olkoon \((x_n)\) lukujono, \(x_n = -n, n \in \NN\). On voimassa
  \(0 \geq x_n\) kaikilla \(n \in \NN\), eli \((x_n)\) on ylhäältä
  rajoitettu.

  Oletetaan, että \(M \in \RR\) on jonon \((x_n)\) alaraja, eli
  \(M \leq x_n\) kaikilla \(n \in \NN\).
  Valitaan \(n \in \NN\) siten, että \(n > -M\). Sitten
  \(x_n = -n < M\). Tämä on ristiriita oletuksen
  \(M \leq x_n\) kaikilla \(n \in \NN\) kanssa. Jonolla \((x_n)\) ei
  siis ole alarajaa.
\end{enumerate}
\end{solution}

\begin{exercise}
Matroskin ja Musti ovat lukeneet Fedja-sedän matikanvihkoa. Molemmat ovat
kopioineet omaan vihkoonsa lukujonon suppenemisen määritelmän. Kopioitaessa
määritelmä on kuitenkin muuttunut. Tutki ovatko Matroskinin ja Mustin
määritelmät ekvivalentteja alkuperäisen määritelmän kanssa.

\textbf{Matroskin:} Lukujono $(x_n)$ suppenee kohti lukua $a$, jos on
olemassa sellainen luku $\lambda$, että kaikilla $m$ pätee
$\abs{x_n-a}<\lambda$ kun $n>m$.

\textbf{Musti:} Lukujono $(x_n)$ suppenee kohti lukua $a$, jos jokaisella
$\xi>0$ on sellainen $\mu$, että $\abs{x_n-a}<7\xi$ kun $n>\mu/8$.
\end{exercise}

\begin{solution}
\end{solution}

\end{document}
